\section{Markowitz: Model Portofolio Modern}\label{appendix-a}

Model Markowitz, atau dikenal sebagai \textbf{Modern Portfolio Theory (MPT)}, diperkenalkan oleh Harry Markowitz pada tahun 1952. Berikut adalah langkah-langkah penurunan rumusnya dari tingkat filosofis hingga menjadi persamaan optimasi biner.

\begin{center}\rule{0.5\linewidth}{0.5pt}\end{center}

\hypertarget{fondasi-filosofis-dan-aksioma-dasar}{%
\subsection{Fondasi Filosofis dan Aksioma Dasar}\label{fondasi-filosofis-dan-aksioma-dasar}}

Penurunan ini bermula dari beberapa postulat tentang perilaku manusia dan pasar:

\hypertarget{a.-aksioma-rasionalitas-risk-aversion}{%
\subsubsection{Aksioma Rasionalitas (Risk Aversion)}\label{a.-aksioma-rasionalitas-risk-aversion}}

Investasi bukan sekadar mencari keuntungan sebesar-besarnya. Investor diasumsikan rasional: - Jika ada dua aset dengan \textbf{return yang sama}, investor akan memilih yang memiliki \textbf{risiko lebih rendah}. - Jika ada dua aset dengan \textbf{risiko yang sama}, investor akan memilih yang memiliki \textbf{return lebih tinggi}.

\hypertarget{b.-konvensi-risiko-sebagai-volatilitas}{%
\subsubsection{Konvensi Risiko sebagai Volatilitas}\label{b.-konvensi-risiko-sebagai-volatilitas}}

Dalam pandangan Markowitz, risiko bukanlah ``kemungkinan bangkrut'', melainkan \textbf{ketidakpastian}. Secara statistik, ketidakpastian diukur sebagai dispersi dari rata-rata, yaitu \textbf{Varians (\(\sigma^2\))} atau \textbf{Deviasi Standar (\(\sigma\))}.

\hypertarget{c.-filosofi-diversifikasi}{%
\subsubsection{Filosofi Diversifikasi}\label{c.-filosofi-diversifikasi}}

``Don't put all your eggs in one basket.'' Markowitz membuktikan secara matematis bahwa menggabungkan aset yang tidak berkorelasi sempurna dapat menurunkan risiko total portofolio tanpa harus mengurangi return yang diharapkan.

\begin{center}\rule{0.5\linewidth}{0.5pt}\end{center}

\hypertarget{postulat-matematis-mean-variance-optimization}{%
\subsection{Postulat Matematis: Mean-Variance Optimization}\label{postulat-matematis-mean-variance-optimization}}

Diberikan \(N\) aset, kita mendefinisikan dua variabel statistik utama: 1. \textbf{Expected Return (\(\mu\)):} Nilai rata-rata keuntungan di masa depan. 2. \textbf{Kovarians (\(\Sigma\)):} Matriks yang merangkum varians tiap aset (diagonal) dan hubungan pergerakan antar aset (off-diagonal).

Tujuan investor adalah meminimalkan risiko total portofolio: \[ \text{Minimize } \sigma_p^2 = \sum_{i,j} w_i w_j \Sigma_{ij} \]

\hypertarget{konstruksi-matriks-kovarians-sigma}{%
\subsubsection{\texorpdfstring{Konstruksi Matriks Kovarians (\(\Sigma\))}{Konstruksi Matriks Kovarians (\textbackslash Sigma)}}\label{konstruksi-matriks-kovarians-sigma}}

Risiko portofolio bukan sekadar jumlah risiko aset individu, melainkan bagaimana aset-aset tersebut bergerak bersama. Matriks \(\Sigma\) (ukuran \(N \times N\)) dirumuskan sebagai:

\[ \Sigma = \begin{pmatrix} 
\sigma_1^2 & \sigma_{12} & \dots & \sigma_{1N} \\
\sigma_{21} & \sigma_2^2 & \dots & \sigma_{2N} \\
\vdots & \vdots & \ddots & \vdots \\
\sigma_{N1} & \sigma_{N2} & \dots & \sigma_N^2
\end{pmatrix} \]

Dimana: 1. \textbf{Elemen Diagonal (\(\sigma_i^2\)):} Merupakan \textbf{Varians} dari aset \(i\). Ini mewakili risiko intrinsik atau volatilitas aset tersebut. 2. \textbf{Elemen Off-Diagonal (\(\sigma_{ij}\)):} Merupakan \textbf{Kovarians} antara aset \(i\) dan \(j\). Kovarians tidak bisa didapatkan hanya dari \(\sigma\) masing-masing aset, melainkan membutuhkan \textbf{Koefisien Korelasi (\(\rho_{ij}\))}.

Hubungan matematisnya adalah: \[ \sigma_{ij} = \rho_{ij} \cdot \sigma_i \cdot \sigma_j \]

Di mana: - \(\sigma_i, \sigma_j\): Deviasi standar (volatilitas) masing-masing aset. - \textbf{Koefisien Korelasi (\(\rho_{ij}\)):} Didapatkan dengan menormalisasi kovarians terhadap produk deviasi standar kedua aset: \[ \rho_{ij} = \frac{\sigma_{ij}}{\sigma_i \cdot \sigma_j} \] - Rentang nilai \(\rho_{ij}\) adalah \([-1, 1]\): - \textbf{1:} Korelasi positif sempurna (bergerak searah). - \textbf{0:} Tidak ada hubungan linear. - \textbf{-1:} Korelasi negatif sempurna (bergerak berlawanan arah).

\hypertarget{bagaimana-sigma-didapatkan}{%
\paragraph{\texorpdfstring{Bagaimana \(\Sigma\) didapatkan?}{Bagaimana \textbackslash Sigma didapatkan?}}\label{bagaimana-sigma-didapatkan}}

Misalkan kita memiliki matriks return terpusat \(\tilde{R}\) berukuran \(T \times N\) (dimana setiap elemen sudah dikurangi rata-ratanya, \(R_{it} - \bar{R}_i\)). Matriks \(\Sigma\) didapat dari perkalian dot product antar kolom aset: \[ \Sigma = \frac{1}{T-1} \tilde{R}^T \tilde{R} \]

Dalam perkalian ini, kolom aset \(i\) yang dikalikan dengan dirinya sendiri (\(\tilde{R}_i \cdot \tilde{R}_i\)) mengisi posisi diagonal sebagai varians, sedangkan perkalian antar kolom berbeda (\(\tilde{R}_i \cdot \tilde{R}_j\)) mengisi posisi off-diagonal sebagai kovarians.

\begin{center}\rule{0.5\linewidth}{0.5pt}\end{center}

\hypertarget{transformasi-ke-model-keputusan-biner-x-in-01}{%
\subsection{\texorpdfstring{Transformasi ke Model Keputusan Biner (\(x \in \{0,1\}\))}{Transformasi ke Model Keputusan Biner (x \textbackslash in \textbackslash\{0,1\textbackslash\})}}\label{transformasi-ke-model-keputusan-biner-x-in-01}}

Dalam model Markowitz asli, variabelnya adalah bobot (\(w \in [0,1]\)). Namun, dalam masalah seleksi aset (seperti pada komputasi kuantum), kita sering menggunakan variabel keputusan biner: - \(x_i = 1\): Aset dipilih untuk masuk portofolio. - \(x_i = 0\): Aset tidak dipilih.

Maka, fungsi objektif dasarnya adalah: \[ \text{Objektif} = \underbrace{x^T \Sigma x}_{\text{Risiko}} - \underbrace{\lambda \mu^T x}_{\text{Return}} \] Dimana \(\lambda\) adalah \textbf{Risk Aversion Coefficient}. Jika \(\lambda\) besar, investor sangat mengejar return. Jika \(\lambda\) kecil, investor sangat menghindari risiko.

\begin{center}\rule{0.5\linewidth}{0.5pt}\end{center}

\hypertarget{penanganan-batasan-the-penalty-term}{%
\subsection{Penanganan Batasan (The Penalty Term)}\label{penanganan-batasan-the-penalty-term}}

Dalam dunia nyata, investasi memiliki batasan (\emph{constraints}), misalnya: ``Saya hanya ingin memilih tepat \(K\) aset dari \(N\) pilihan.'' Secara matematis: \[ \sum_{i=1}^N x_i = K \]

Dalam teknik optimasi konvensional, kita menggunakan metode \textbf{Appendix C}. Namun, agar persamaan ini bisa diselesaikan oleh sistem Ising atau QUBO (yang bersifat \emph{unconstrained}), kita harus mengubah ``batasan keras'' (\emph{hard constraint}) menjadi ``penalti lunak'' (\emph{soft penalty}).

\hypertarget{mengapa-kuadratik}{%
\subsubsection{Mengapa Kuadratik?}\label{mengapa-kuadratik}}

Kita menggunakan bentuk selisih kuadrat: \[ \text{Penalti} = A \left( \sum x_i - K \right)^2 \] 
\begin{itemize}
    \item Jika \(\sum x_i = K\), maka \((K-K)^2 = 0\). Tidak ada penalti (energi minimum).
    \item Jika \(\sum x_i \neq K\), maka nilai tersebut akan dikuadratkan (selalu positif) dan dikalikan dengan konstanta \(A\) yang besar
\end{itemize} 
Hal ini akan meningkatkan total energi secara drastis, sehingga sistem akan dipaksa menjauhi solusi yang melanggar batasan ini.

\begin{center}\rule{0.5\linewidth}{0.5pt}\end{center}

\hypertarget{formulasi-akhir-the-markowitz-lagrangian}{%
\subsection{Formulasi Akhir (The Markowitz Lagrangian)}\label{formulasi-akhir-the-markowitz-lagrangian}}

Dengan menggabungkan semua postulat di atas, kita mendapatkan fungsi energi total (Hamiltonian) yang digunakan dalam optimasi:

\[ \min_{x \in \{0,1\}^N} \mathcal{L}(x) = \underbrace{x^T \Sigma x}_{\text{Postulat Risiko}} - \underbrace{\lambda \mu^T x}_{\text{Aksioma Return}} + \underbrace{A \left(\sum_{i=1}^N x_i - K \right)^2}_{\text{Konvensi Batasan}} \]

Persamaan ini merangkum seluruh filosofi Markowitz: 1. \textbf{Suku pertama} menjaga agar aset yang dipilih tidak saling bergejolak (diversifikasi). 2. \textbf{Suku kedua} menarik sistem ke arah keuntungan yang lebih tinggi. 3. \textbf{Suku ketiga} memastikan solusi yang ditemukan adalah solusi yang valid sesuai batasan jumlah aset.
